\section*{第 \ref{ch:b2:energy-geology} 章——能源地质学}
% numbers from the Book 4 second-writer scratch calculator calc/ch30.py

\begin{solution}{exo:b2:energy-geology:1}
(a) 1.58：I 型；0.95：III 型；1.24：II 型；1.12：II/III 混合型。III 型干酪根在
 0.95 时主要生成天然气。(b) H/C $= 150/100 = 1.50$，O/C $= 8/100 = 0.08$。(c) 以 \qty{100}{g} 计：
$77.4/12.011 = \qty{6.444}{mol}$ C, $9.8/1.008 = \qty{9.722}{mol}$ H, $8.2/15.999 =
\qty{0.513}{mol}$ O；H/C $= 1.51$ 和 O/C $= 0.0795$，均与分子式相符。
\end{solution}

\begin{solution}{exo:b2:energy-geology:2}
(a) 0.45：未成熟，低于 \qty{0.60}{\percent} 的起始值；0.68：正在生油；1.1：仍在
生油窗内，接近其末端；1.6：超过 1.20--\qty{1.30}{\percent}，油裂解为甲烷，
岩石生成天然气；2.5：仅生成天然气，生成作用持续至约 \qty{4.8}{\percent}。(b) 否：III 型干酪根的甲烷
生成始于约 \qty{0.73}{\percent}。(c) $F = (\ln R_o + 1.6)/3.7$：$(-0.511 +
1.6)/3.7 = 0.294$ 和 $(0.262 + 1.6)/3.7 = 0.503$。
\end{solution}

\begin{solution}{exo:b2:energy-geology:3}
(a) 储层中，$6 \times 10^6 \times 25 \times 0.18 \times 0.70 = \qty{1.89e7}{m^3}$；
$\mathrm{STOIIP} = 1.89 \times 10^7/1.25 = \qty{1.51e7}{m^3}$. (b) $0.30 \times 1.512 \times 10^7 =
\qty{4.54e6}{m^3}$ 和 $0.40 \times 1.512 \times 10^7 = \qty{6.05e6}{m^3}$。
\end{solution}

\begin{solution}{exo:b2:energy-geology:4}
(a) $V = 10^7 \times 200 = \qty{2e9}{m^3}$ 和 $q_R = 2.5 \times 10^6 \times 2 \times 10^9 \times 60
= \qty{3.0e17}{J}$。(b) 低于 \qty{90}{\celsius} 时属\omterm{def:b2:energy-geology:geothermal}{低焓资源}，供直接利用
热量。(c) $P = 0.25 \times 3.0 \times 10^{17}/(30 \times \qty{3.156e7}{s}) = \qty{79.2}{MW}$ 的热功率。
\end{solution}

\begin{solution}{exo:b2:energy-geology:5}
(a) $k_2/k_1 = \exp[(E/R)(1/T_1 - 1/T_2)] = 2$，代入 $1/373.15 - 1/383.15 = \qty{6.994e-5}{K^{-1}}$
得 $E = 8.314 \times \ln 2/(6.994 \times 10^{-5}) = \qty{82.4}{kJ/mol}$。(b) $\exp(26800 \times
6.994 \times 10^{-5}) = 6.52$。(c) 倍增规则的有效能较低，相对而言给予
低温下耗时较长的阶段更大的权重；对于具有较高
活化能的 I 型干酪根，TTI 会过早地预报生烃。
\end{solution}

\begin{solution}{exo:b2:energy-geology:6}
(a) 近似下，$x = 1/2$ 给出 $(ART^2/\beta E)\,e^{-E/RT} = \ln 2$；取对数，
$E/RT = \ln[AT^2/(\beta(E/R)\ln 2)]$，于是得出所述形式。从任一接近
\qty{400}{K} 的猜测出发迭代，$\beta$ 以开尔文每秒计；在 \qty{1}{\celsius\per\Ma}
（$\beta = \qty{3.169e-14}{K/s}$），$T = \qty{409.6}{K}$，\qty{136.4}{\celsius}；在每分钟 \qty{10}{K}
的加热速率下（$\beta = \qty{0.1667}{K/s}$，速率快 $\num{5.26e12}$ 倍），$T = \qty{719.0}{K}$，
\qty{445.9}{\celsius}。(b) 差值为 \qty{309.5}{K}；估计值 $6.259 \times
\ln(5.26 \times 10^{12}) = 6.259 \times 29.29 = \qty{183}{K}$ 远偏小，因为 $RT^2/E$
本身随 $T$ 增大（在 \qty{719}{K} 时达 \qty{19.3}{K}）：线性位移仅适用于
速率变化较小的情形。
\end{solution}

\begin{solution}{exo:b2:energy-geology:7}
$\lambda_5 - \lambda_8 = \qty{8.29725e-10}{\text{年}^{-1}}$ 和 $r_0 = 0.7204/99.2742 = 0.0072567$。
(a) 在 \qty{2.0}{\Ga} 时，$r = 0.0072567\,e^{1.6595} = 0.03814$，份额 $0.03814/1.03814 =
\qty{3.67}{\percent}$；在 \qty{1.8}{\Ga} 时，$r = 0.03231$，份额 \qty{3.13}{\percent}。(b) $r =
3/97 = 0.030928$ 和 $t = \ln(0.030928/0.0072567)/(8.29725 \times 10^{-10}) = \qty{1.75}{\Ga}$。
(c) 贫化矿石与天然铀自那时起都按相同倍数衰变，故二者的
$^{235}$U 份额之比在今天与反应刚结束时相同：矿石当时的份额为
正常份额的一半，因此约一半的 $^{235}$U 已发生裂变。
\end{solution}

\begin{solution}{exo:b2:energy-geology:8}
(a) $10^9/1.872 = \qty{5.34e8}{m^3}$；$10^9/621.4 = \qty{1.61e6}{m^3}$；$10^9/686.4 =
\qty{1.46e6}{m^3}$。(b) $7 \times 10^9/583 = \qty{1.20e7}{m^3}$ 的 \ce{CO2}，充填
$0.35 \times 0.30 = 0.105$ 的岩石体积：$1.20 \times 10^7/0.105 = \qty{1.14e8}{m^3}$ 的岩石，
平均厚度为 $1.143 \times 10^8/(5 \times 10^6) = \qty{22.9}{m}$。
\end{solution}

\begin{solution}{exo:b2:energy-geology:9}
(a) $V = \qty{9e9}{m^3}$ 和 $q_R = 2.6 \times 10^6 \times 9 \times 10^9 \times 235 =
\qty{5.50e18}{J}$。(b) 在 $30 \times \qty{3.156e7}{s} = \qty{9.47e8}{s}$ 内：$R_g = 0.1$ 时给出
\qty{581}{MW}，$0.25$ 时给出 \qty{1450}{MW} 的热功率。(c) $6 \times (10 \text{ 至 } 20) = 60$ 至
\qty{120}{MW} 的电功率，$60/581 = 0.10$ 至 $120/581 = 0.21$ 占热功率的比例：大部分
热量被排放，与任何在较小温差下工作的电厂一样。
\end{solution}

\begin{solution}{exo:b2:energy-geology:10}
(a) 强度为 $0.6 \times (80 - 39) = \qty{24.6}{MPa} > \qty{22}{MPa}$：稳定。破裂需要
$0.6\,\Delta p = 24.6 - 22$，$\Delta p = \qty{4.33}{MPa} = \qty{43}{bar}$，与
限定丹佛地震的 \qty{32}{bar} 同量级。(b) $(24.6 - 24)/0.6 = \qty{1.0}{MPa} =
\qty{10}{bar}$：已接近破裂的断层只需较小的升压。(c) 注入的流体提高
了一带的压力，该带借扩散通过透水岩石和断层带，历经数年向数千米外扩展；
只要遇到已接近破裂的断层，较小的升压便已足够。
\end{solution}

\begin{solution}{exo:b2:energy-geology:11}
(a) 记 $x = e^{-u}$，其中 $u = (ART^2/\beta E)\,e^{-E/RT}$。转化率为
\qty{10}{\percent} 即 $x = 0.9$，$u = \ln(10/9)$；\qty{90}{\percent} 即 $u = \ln 10$。由于
$\mathrm{d}\ln u/\mathrm{d}T = E/RT^2 + 2/T \approx E/RT^2$, $\Delta(\ln u) \approx (E/RT^2)\,\Delta T$,
由此得出公式，其中 $\ln\ln 10 - \ln\ln(10/9) = 0.834 + 2.250 = 3.084$。(b) 取 $T =
\qty{409.6}{K}$，$RT^2/E = \qty{6.26}{K}$ 和 $\Delta T = \qty{19.3}{K}$；精确的单一
活化能曲线给出 \qty{18.6}{K}，自 125.3 升至 \qty{143.8}{\celsius}。(c) 实线曲线自
\qty{10}{\percent}（约 \qty{122}{\celsius}）升至 \qty{90}{\percent}（约 \qty{148}{\celsius}），在约
\qty{26}{K} 的区间内完成，活化能分布的展宽拓宽了转化区间。
\end{solution}

\begin{solution}{exo:b2:energy-geology:12}
(a) $2 \times 10^{10}/20 = \qty{1.0e9}{m^3}$（\qty{1.0}{km^3}）和 $2 \times 10^{10}/60 =
\qty{3.3e8}{m^3}$。(b) $2 \times 10^{10}/621.4 = \qty{3.22e7}{m^3}$。(c) $3.22 \times 10^7/0.02 =
\qty{1.61e9}{m^3}$ 的孔隙，若每平方米孔隙空间为 $50 \times 0.30 = \qty{15}{m}$，
则需要面积 $1.61 \times 10^9/15 = \qty{1.07e8}{m^2}$，约 \qty{107}{km^2}。
\end{solution}

\begin{solution}{pb:b2:energy-geology:1}
\noindent\textbf{第一部分。}
\textbf{1.} 在 \qty{70}{\Ma} 时：$50 \times 60 = \qty{3000}{m}$，$15 + 25 \times 3.0 =
\qty{90}{\celsius}$。今天：$3000 + 70 \times 30 = \qty{5100}{m}$，\qty{142.5}{\celsius}。
\textbf{2.} $60 \times 0.025 = \qty{1.5}{\celsius\per\Ma}$，然后 $30 \times 0.025 =
\qty{0.75}{\celsius\per\Ma}$。
\textbf{3.} 15--\qty{20}{\celsius}：\qty{3.33}{\Ma}；自 20 至 \qty{90}{\celsius} 的每段：
$10/1.5 = \qty{6.67}{\Ma}$；自 90 至 \qty{140}{\celsius} 的每段：$10/0.75 = \qty{13.3}{\Ma}$；
140--\qty{142.5}{\celsius}：\qty{3.33}{\Ma}。
\textbf{4.} 各区间的值 $\Delta t_n 2^n$，并给出累计和（括号中）：15--20（$n = -9$）0.0065；20--30
0.026 (0.033); 30--40 0.052 (0.085); 40--50 0.104 (0.189); 50--60 0.208 (0.397); 60--70 0.417
(0.814); 70--80 0.833 (1.65); 80--90 1.67 (3.31); 90--100 6.67 (9.98); 100--110 13.3 (23.3);
110--120 26.7 (50.0); 120--130 53.3 (103); 130--140 107 (210); 140--142.5 ($n = 4$) 53.3：
TTI 今天为 263。
\textbf{5.} 15 在 100--\qty{110}{\celsius} 区间被越过，其温度为 $100 + 10 \times
(15 - 9.98)/13.33 = \qty{103.8}{\celsius}$，在 \qty{70}{\Ma} 后 $(103.8 - 90)/0.75 = \qty{18.4}{\Ma}$ 时达到
即 \qty{51.6}{\Ma} 时，深度为 $(103.8 - 15)/0.025 = \qty{3550}{m}$。
\textbf{6.} 75：\qty{124.7}{\celsius}，\qty{23.7}{\Ma}，\qty{4390}{m}。160：\qty{135.3}{\celsius}，
\qty{9.6}{\Ma}，\qty{4810}{m}。
\textbf{7.} 否：TTI 为 263 时已超过生油终点（160），约自 \qty{10}{\Ma} 前起。

\noindent\textbf{第二部分。}
\textbf{8.} $k = 9.21 \times 10^{13}\,e^{-26800/363.15} = \qty{8.20e-19}{s^{-1}}$。该岩石
一直低于 \qty{90}{\celsius}，在 \qty{70}{\Ma} 之前的 \qty{50}{\Ma} 内，故转化分数
小于 $k\,t = 8.20 \times 10^{-19} \times 50 \times 3.156 \times 10^{13} = \num{1.29e-3}$。
\textbf{9.} $\beta = 0.75/(3.156 \times 10^{13}) = \qty{2.376e-14}{K/s}$, $T = \qty{415.65}{K}$:
$ART^2/(\beta E) = \num{2.498e28}$ 与 $e^{-E/RT} = \num{9.95e-29}$，故 $u = 2.486$，$x = 0.0833$ 且
转化率为 0.917。
\textbf{10.} \qty{134.7}{\celsius}.
\textbf{11.} $70 - (134.7 - 90)/0.75 = \qty{10.4}{\Ma}$.
\textbf{12.} TTI 给出生油峰值出现在 \qty{23.7}{\Ma}，动力学结果则晚 \qty{13}{\Ma}，
原因在于倍增规则对应 $E \approx \qty{82}{kJ/mol}$，且对低温段赋予了
过大的权重，对活化能为 \qty{223}{kJ/mol} 的干酪根尤为不当（见练习~\ref{exo:b2:energy-geology:5}）。
\textbf{13.} 更早：较低的活化能使半转化温度降低。

\noindent\textbf{第三部分。}
\textbf{14.} 岩石：$4 \times 10^8 \times 50 \times 2400 = \qty{4.8e13}{kg}$；潜力 $0.030
\times 4.8 \times 10^{13} = \qty{1.44e12}{kg}$；生成量 $0.92 \times 1.44 \times 10^{12} =
\qty{1.32e12}{kg}$.
\textbf{15.} $1.3248 \times 10^{12}/850 = \qty{1.56e9}{m^3}$.
\textbf{16.} $1.2 \times 10^7 \times 40 \times 0.22 \times 0.75/1.30 = \qty{6.09e7}{m^3}$.
\textbf{17.} $6.09 \times 10^7/1.56 \times 10^9 = 0.039$，\qty{3.9}{\percent}：大部分生成的油
仍留在烃源岩中，沿运移路径散失或逸散至地表。
\textbf{18.} $0.35 \times 6.09 \times 10^7 = \qty{2.13e7}{m^3}$.
\textbf{19.} 一次和二次采收后，约三分之二的油仍留在储层中；
\ce{CO2} 注入可进一步采出\omterm{def:b2:palaeoenvironments:assemblages}{原地}油量的 10 至 \qty{15}{\percent}。

\noindent\textbf{第四部分。}
\textbf{20.} $15 + 25 \times 2 = \qty{65}{\celsius}$，且 $0.101 + 1000 \times 9.807 \times 2000/10^6
= \qty{19.7}{MPa}$：两者均高于临界点（\qty{31.0}{\celsius}，\qty{7.38}{MPa}），因此
\ce{CO2} 处于超临界状态，密度约为 \qty{686}{kg.m^{-3}}。
\textbf{21.} $2.13 \times 10^7 \times 1.30 = \qty{2.77e7}{m^3}$.
\textbf{22.} 孔隙体积 $1.2 \times 10^7 \times 40 \times 0.22 = \qty{1.056e8}{m^3}$，且 $E =
2.77 \times 10^7/1.056 \times 10^8 = 0.262$，远高于盐水地层的 1 至 \qty{4}{\percent}：
采出的油已掏空了一个已知的封闭空间。
\textbf{23.} $6 \times 10^8 \times 80 \times 0.20 \times 691 \times 0.02 = \qty{1.33e11}{kg}$,
\qty{133}{Mt}.
\textbf{24.} 按每年 \qty{1}{Mt} 计，约 133 年。
\textbf{25.} $M = \rho\,N_p\,B_o = 686 \times 2.77 \times 10^7 = \qty{1.90e10}{kg}$：枯竭的
储层可储存 \textbf{\qty[text-series-to-math=true]{19.0}{Mt}} 的 \ce{CO2}。
\end{solution}
